%% ============================================================
%% nro/competencias.tex: as competências do egresso, pela DCN do curso
%%
%% Mantido pela DIRETORIA. O curso da pessoa (estudante/curso, que o
%% SOMA preenche pela lista dos cursos) escolhe a lista; a matriz da
%% Parte IV e a ficha de cada atividade usam a dela. Curso fora das
%% listas usa a "geral". A atividade cita as suas em
%% competencias = {I, III, V}.
%%
%% Cursos com a mesma base de diretrizes dividem a lista. As de
%% engenharia são as competências gerais do art. 4º da Resolução CNE/CES
%% nº 2/2019, mais as específicas da área do curso (o parágrafo único do
%% mesmo artigo deixa ao curso as específicas).
%%
%% O SOMA oferece os mesmos cursos (mod-aproveitamento.js, AAC_CURSOS):
%% curso novo entra aqui e lá.
%% ============================================================

%% ------------------------------------------------------------
%% Engenharia: circuitos, energia, controle, sistemas e computação
%% ------------------------------------------------------------
\conjuntocompetencias{eng-eletrica}{o art.~4\textordmasculine\ da Resolução CNE/CES n\textordmasculine~2/2019 (Diretrizes Curriculares Nacionais dos cursos de Engenharia) e as competências específicas do curso}
\competencia{I}{Formular e conceber soluções desejáveis de engenharia, analisando e compreendendo os usuários dessas soluções e o seu contexto.}
\competencia{II}{Analisar e compreender os fenômenos físicos e químicos por meio de modelos simbólicos, físicos e outros, verificados e validados por experimentação.}
\competencia{III}{Conceber, projetar e analisar sistemas, produtos (bens e serviços), componentes ou processos.}
\competencia{IV}{Implantar, supervisionar e controlar as soluções de engenharia.}
\competencia{V}{Comunicar-se eficazmente nas formas escrita, oral e gráfica.}
\competencia{VI}{Trabalhar e liderar equipes multidisciplinares.}
\competencia{VII}{Conhecer e aplicar com ética a legislação e os atos normativos no âmbito do exercício da profissão.}
\competencia{VIII}{Aprender de forma autônoma e lidar com situações e contextos complexos, atualizando-se em relação aos avanços da ciência, da tecnologia e aos desafios da inovação.}
\competencia{IX}{Projetar, implementar e ensaiar circuitos e sistemas eletrônicos, de potência e de instrumentação, do esquemático à placa.}
\competencia{X}{Modelar, simular e controlar sistemas dinâmicos e sinais, e integrar hardware, firmware e software em sistemas embarcados.}
\cursos{eng-eletrica}{Engenharia Elétrica; Engenharia de Controle e Automação; Engenharia de Sistemas; Engenharia de Computação}

%% ------------------------------------------------------------
%% Engenharia: projeto mecânico, fabricação e produção
%% ------------------------------------------------------------
\conjuntocompetencias{eng-mecanica}{o art.~4\textordmasculine\ da Resolução CNE/CES n\textordmasculine~2/2019 (Diretrizes Curriculares Nacionais dos cursos de Engenharia) e as competências específicas do curso}
\competencia{I}{Formular e conceber soluções desejáveis de engenharia, analisando e compreendendo os usuários dessas soluções e o seu contexto.}
\competencia{II}{Analisar e compreender os fenômenos físicos e químicos por meio de modelos simbólicos, físicos e outros, verificados e validados por experimentação.}
\competencia{III}{Conceber, projetar e analisar sistemas, produtos (bens e serviços), componentes ou processos.}
\competencia{IV}{Implantar, supervisionar e controlar as soluções de engenharia.}
\competencia{V}{Comunicar-se eficazmente nas formas escrita, oral e gráfica.}
\competencia{VI}{Trabalhar e liderar equipes multidisciplinares.}
\competencia{VII}{Conhecer e aplicar com ética a legislação e os atos normativos no âmbito do exercício da profissão.}
\competencia{VIII}{Aprender de forma autônoma e lidar com situações e contextos complexos, atualizando-se em relação aos avanços da ciência, da tecnologia e aos desafios da inovação.}
\competencia{IX}{Projetar, dimensionar e especificar componentes e sistemas mecânicos, com a seleção de materiais e de processos de fabricação.}
\competencia{X}{Analisar e simular o comportamento estrutural, térmico e de escoamento de sistemas, e planejar a produção, a qualidade e a manutenção.}
\cursos{eng-mecanica}{Engenharia Mecânica; Engenharia Aeroespacial; Engenharia de Produção}

%% ------------------------------------------------------------
%% Engenharia: processos e materiais
%% ------------------------------------------------------------
\conjuntocompetencias{eng-processos}{o art.~4\textordmasculine\ da Resolução CNE/CES n\textordmasculine~2/2019 (Diretrizes Curriculares Nacionais dos cursos de Engenharia) e as competências específicas do curso}
\competencia{I}{Formular e conceber soluções desejáveis de engenharia, analisando e compreendendo os usuários dessas soluções e o seu contexto.}
\competencia{II}{Analisar e compreender os fenômenos físicos e químicos por meio de modelos simbólicos, físicos e outros, verificados e validados por experimentação.}
\competencia{III}{Conceber, projetar e analisar sistemas, produtos (bens e serviços), componentes ou processos.}
\competencia{IV}{Implantar, supervisionar e controlar as soluções de engenharia.}
\competencia{V}{Comunicar-se eficazmente nas formas escrita, oral e gráfica.}
\competencia{VI}{Trabalhar e liderar equipes multidisciplinares.}
\competencia{VII}{Conhecer e aplicar com ética a legislação e os atos normativos no âmbito do exercício da profissão.}
\competencia{VIII}{Aprender de forma autônoma e lidar com situações e contextos complexos, atualizando-se em relação aos avanços da ciência, da tecnologia e aos desafios da inovação.}
\competencia{IX}{Projetar, operar e otimizar processos de transformação física e química, com balanços de massa e de energia.}
\competencia{X}{Caracterizar, selecionar e processar materiais e insumos, com controle de qualidade e de segurança.}
\cursos{eng-processos}{Engenharia Química; Engenharia Metalúrgica; Engenharia de Materiais; Engenharia de Minas; Engenharia de Alimentos}

%% ------------------------------------------------------------
%% Engenharia: infraestrutura e ambiente
%% ------------------------------------------------------------
\conjuntocompetencias{eng-civil}{o art.~4\textordmasculine\ da Resolução CNE/CES n\textordmasculine~2/2019 (Diretrizes Curriculares Nacionais dos cursos de Engenharia) e as competências específicas do curso}
\competencia{I}{Formular e conceber soluções desejáveis de engenharia, analisando e compreendendo os usuários dessas soluções e o seu contexto.}
\competencia{II}{Analisar e compreender os fenômenos físicos e químicos por meio de modelos simbólicos, físicos e outros, verificados e validados por experimentação.}
\competencia{III}{Conceber, projetar e analisar sistemas, produtos (bens e serviços), componentes ou processos.}
\competencia{IV}{Implantar, supervisionar e controlar as soluções de engenharia.}
\competencia{V}{Comunicar-se eficazmente nas formas escrita, oral e gráfica.}
\competencia{VI}{Trabalhar e liderar equipes multidisciplinares.}
\competencia{VII}{Conhecer e aplicar com ética a legislação e os atos normativos no âmbito do exercício da profissão.}
\competencia{VIII}{Aprender de forma autônoma e lidar com situações e contextos complexos, atualizando-se em relação aos avanços da ciência, da tecnologia e aos desafios da inovação.}
\competencia{IX}{Projetar, dimensionar e analisar estruturas, obras e sistemas de infraestrutura.}
\competencia{X}{Avaliar impactos ambientais e propor soluções de saneamento, de uso do solo e de gestão de recursos naturais.}
\cursos{eng-civil}{Engenharia Civil; Engenharia Ambiental; Engenharia Agrícola e Ambiental; Engenharia Florestal}

%% ------------------------------------------------------------
%% Saúde: movimento, funcionalidade e reabilitação
%% ------------------------------------------------------------
\conjuntocompetencias{saude-reabilitacao}{as Diretrizes Curriculares Nacionais do curso (as competências gerais e as específicas)}
\competencia{I}{Atenção à saúde: desenvolver ações de prevenção, promoção, proteção e reabilitação da saúde, individual e coletiva, com qualidade e segundo os princípios da ética.}
\competencia{II}{Tomada de decisões: avaliar, sistematizar e decidir as condutas mais adequadas, com base em evidências científicas.}
\competencia{III}{Comunicação: comunicar-se com usuários, profissionais e o público, nas formas verbal, não verbal e escrita, com confidencialidade.}
\competencia{IV}{Liderança: trabalhar em equipe multiprofissional e assumir posições de liderança, com compromisso, responsabilidade e empatia.}
\competencia{V}{Administração e gerenciamento: gerenciar recursos físicos, materiais, tecnológicos e de informação, e empreender.}
\competencia{VI}{Educação permanente: aprender continuamente, na formação e na prática, e contribuir para a formação de outros profissionais.}
\competencia{VII}{Avaliar o movimento humano e a funcionalidade, e intervir com recursos terapêuticos e com tecnologias assistivas e de reabilitação.}
\competencia{VIII}{Desenvolver, adaptar e avaliar dispositivos e tecnologias para a reabilitação, a participação e a autonomia de pessoas com deficiência.}
\cursos{saude-reabilitacao}{Fisioterapia; Terapia Ocupacional; Educação Física}

%% ------------------------------------------------------------
%% Saúde: cuidado clínico
%% ------------------------------------------------------------
\conjuntocompetencias{saude-clinica}{as Diretrizes Curriculares Nacionais do curso (as competências gerais e as específicas)}
\competencia{I}{Atenção à saúde: desenvolver ações de prevenção, promoção, proteção e reabilitação da saúde, individual e coletiva, com qualidade e segundo os princípios da ética.}
\competencia{II}{Tomada de decisões: avaliar, sistematizar e decidir as condutas mais adequadas, com base em evidências científicas.}
\competencia{III}{Comunicação: comunicar-se com usuários, profissionais e o público, nas formas verbal, não verbal e escrita, com confidencialidade.}
\competencia{IV}{Liderança: trabalhar em equipe multiprofissional e assumir posições de liderança, com compromisso, responsabilidade e empatia.}
\competencia{V}{Administração e gerenciamento: gerenciar recursos físicos, materiais, tecnológicos e de informação, e empreender.}
\competencia{VI}{Educação permanente: aprender continuamente, na formação e na prática, e contribuir para a formação de outros profissionais.}
\competencia{VII}{Aplicar evidências científicas e tecnologias em saúde ao diagnóstico, ao tratamento e ao acompanhamento das pessoas.}
\competencia{VIII}{Participar da pesquisa e da inovação em saúde, com rigor metodológico e respeito às normas éticas de pesquisa com seres humanos.}
\cursos{saude-clinica}{Medicina; Enfermagem; Nutrição; Fonoaudiologia; Tecnologia em Radiologia}

%% ------------------------------------------------------------
%% Saúde: gestão de serviços
%% ------------------------------------------------------------
\conjuntocompetencias{saude-gestao}{as Diretrizes Curriculares Nacionais do curso (as competências gerais e as específicas)}
\competencia{I}{Atenção à saúde: desenvolver ações de prevenção, promoção e proteção da saúde, individual e coletiva, com qualidade e segundo os princípios da ética.}
\competencia{II}{Tomada de decisões: avaliar, sistematizar e decidir com base em evidências e em indicadores.}
\competencia{III}{Comunicação: comunicar-se com usuários, profissionais e o público, nas formas verbal, não verbal e escrita, com confidencialidade.}
\competencia{IV}{Liderança: trabalhar em equipe multiprofissional e assumir posições de liderança, com compromisso, responsabilidade e empatia.}
\competencia{V}{Administração e gerenciamento: gerenciar recursos físicos, materiais, tecnológicos e de informação, e empreender.}
\competencia{VI}{Educação permanente: aprender continuamente, na formação e na prática, e contribuir para a formação de outros profissionais.}
\competencia{VII}{Planejar, organizar e avaliar serviços, processos e tecnologias em saúde, com base em indicadores.}
\competencia{VIII}{Analisar a incorporação de tecnologias em saúde quanto à efetividade, ao custo e ao acesso.}
\cursos{saude-gestao}{Gestão de Serviços de Saúde}

%% ------------------------------------------------------------
%% Computação
%% ------------------------------------------------------------
\conjuntocompetencias{computacao}{a Resolução CNE/CES n\textordmasculine~5/2016 (Diretrizes Curriculares Nacionais dos cursos de Computação)}
\competencia{I}{Identificar problemas com solução algorítmica e conceber soluções computacionais, considerando os limites da computação.}
\competencia{II}{Resolver problemas com ambientes de programação, estruturas de dados e algoritmos adequados.}
\competencia{III}{Projetar, implementar e avaliar sistemas computacionais, com conhecimento do hardware e da infraestrutura de software.}
\competencia{IV}{Compreender e explicar as dimensões quantitativas de um problema, com modelos e análise de dados.}
\competencia{V}{Preparar e apresentar trabalhos técnicos e as suas soluções a públicos diversos, nas formas oral e escrita.}
\competencia{VI}{Trabalhar em equipe e exercer liderança e coordenação em projetos de computação.}
\competencia{VII}{Agir com ética, considerando os impactos sociais, legais e ambientais das soluções computacionais.}
\competencia{VIII}{Gerir a própria aprendizagem e adaptar-se a novas tecnologias e ambientes de trabalho.}
\cursos{computacao}{Ciência da Computação; Sistemas de Informação; Ciência de Dados; Matemática Computacional}

%% ------------------------------------------------------------
%% Ciências exatas
%% ------------------------------------------------------------
\conjuntocompetencias{exatas}{as Diretrizes Curriculares Nacionais do curso}
\competencia{I}{Formular problemas e resolvê-los com modelos matemáticos, físicos ou químicos, com rigor e senso crítico.}
\competencia{II}{Planejar e realizar experimentos e medições, e analisar os dados com a quantificação da incerteza.}
\competencia{III}{Usar métodos numéricos, estatísticos e ferramentas computacionais na análise e na simulação.}
\competencia{IV}{Aplicar o conhecimento da área a problemas tecnológicos e interdisciplinares.}
\competencia{V}{Comunicar resultados científicos com clareza, nas formas oral e escrita.}
\competencia{VI}{Trabalhar em equipes multidisciplinares.}
\competencia{VII}{Agir com ética e com responsabilidade social e ambiental.}
\competencia{VIII}{Aprender de forma autônoma e atualizar-se em relação aos avanços da ciência.}
\cursos{exatas}{Matemática; Estatística; Ciências Atuariais; Física; Química; Química Tecnológica}

%% ------------------------------------------------------------
%% Gestão, contabilidade e economia
%% ------------------------------------------------------------
\conjuntocompetencias{gestao}{as Diretrizes Curriculares Nacionais do curso}
\competencia{I}{Analisar problemas organizacionais e econômicos e propor soluções com base em dados e em evidências.}
\competencia{II}{Planejar, organizar e controlar projetos, processos e recursos.}
\competencia{III}{Avaliar a viabilidade econômica e financeira de projetos e de inovações.}
\competencia{IV}{Gerir pessoas e equipes, e exercer liderança.}
\competencia{V}{Comunicar-se e negociar com clareza, nas formas oral e escrita.}
\competencia{VI}{Usar ferramentas quantitativas e tecnologias de informação na tomada de decisão.}
\competencia{VII}{Agir com ética, responsabilidade social e ambiental e respeito à legislação.}
\competencia{VIII}{Aprender de forma autônoma e lidar com contextos complexos e com a inovação.}
\cursos{gestao}{Administração; Ciências Contábeis; Ciências Econômicas; Controladoria e Finanças; Relações Econômicas Internacionais}

%% ------------------------------------------------------------
%% A lista geral: curso que não está acima
%% ------------------------------------------------------------
\conjuntocompetencias{geral}{as Diretrizes Curriculares Nacionais do curso}
\competencia{I}{Aplicar os conhecimentos do curso a problemas reais, do levantamento de requisitos à validação.}
\competencia{II}{Conceber, projetar e avaliar soluções, com método e com base em evidências.}
\competencia{III}{Planejar e realizar ensaios, medições ou estudos, e analisar os resultados.}
\competencia{IV}{Comunicar-se nas formas escrita, oral e gráfica.}
\competencia{V}{Trabalhar e liderar equipes multidisciplinares.}
\competencia{VI}{Agir com ética e conhecer as normas aplicáveis à atividade.}
\competencia{VII}{Aprender de forma autônoma e lidar com situações complexas.}
